% GENERATED FILE. Edit canonical lessons, then run npm run generate:books.
% canonical-source-hash: fnv1a64:f1ab1cb6406d834b
% canonical-lessons: SA-C344-pratipaksah, SA-C344-virodhah, SA-C344-samarthayati, SA-C344-anumanyate, SA-C344-nanu

\chapter{Taking Sides}
\label{ch:sa-taking-sides}
\hlchaptermodality{344}

\begin{chapteropening}
\textbf{By the end of this chapter:} \emph{I can say who supports a view and who opposes it.}

Complete the last lesson of chapter 344: I can say who supports a view and who opposes it.
\end{chapteropening}

\section[\texorpdfstring{\sk{प्रतिपक्षः} (pratipakṣaḥ)}{प्रतिपक्षः (pratipakṣaḥ)}]{\sk{प्रतिपक्षः} (pratipakṣaḥ) — the opposing side}

\label{lesson:SA-C344-pratipaksah}



\begin{quote}
Before the new one: say the Sanskrit for a listener, then the Sanskrit for sets out (a case).
\end{quote}

\subsection*{You'll want to know: \sk{प्रतिपक्षः}}
\textbf{\sk{प्रतिपक्षः}} — \emph{pratipakṣaḥ} — \textquotedblleft{}the opposing side\textquotedblright{}.

\begin{quote}
\sk{प्रतिपक्षस्य} \sk{मतम्} \sk{अन्यत्} \sk{अस्ति।}
\end{quote}

\begin{quote}
\emph{pratipakṣasya matam anyat asti.}
\end{quote}

\begin{quote}
The other side's view is different.
\end{quote}

A masculine noun.

\begin{cousinweb}
\textbf{\sk{प्रतिपक्षः}} (\emph{pratipakṣaḥ}) is built from \emph{prati}, \textquotedblleft{}against\textquotedblright{}, and \emph{pakṣa}, \textquotedblleft{}a wing, a side\textquotedblright{}; MW gives \emph{prati-pakṣa} \textquotedblleft{}the opposite side, a hostile party\textquotedblright{} (MW).
\end{cousinweb}

\begin{grammarlens}[title={What you've built}]
One more word for giving reasons and opinions.
\end{grammarlens}

\subsection*{Your turn}
\begin{itemize}
\raggedright
  \item \emph{Say it:} \emph{pratipakṣaḥ}
  \item \emph{Say it:} \emph{pratipakṣaḥ}, once more
  \item \emph{Say it:} \emph{pratipakṣasya matam anyat asti.} — The other side's view is different.
  \item \emph{Recall:} say the Sanskrit for a listener, then the Sanskrit for sets out (a case), then say \emph{pratipakṣaḥ} again
\end{itemize}

\begin{tcolorbox}[breakable,colback=teal!4,colframe=teal!35!black,title={Before you move on}]
What does \sk{प्रतिपक्षः} mean? (\textquotedblleft{}The opposing side\textquotedblright{}.) Say it once more.
\end{tcolorbox}
\section[\texorpdfstring{\sk{विरोधः} (virodhaḥ)}{विरोधः (virodhaḥ)}]{\sk{विरोधः} (virodhaḥ) — opposition}

\label{lesson:SA-C344-virodhah}



\begin{quote}
Before the new one: say the Sanskrit for sets out (a case), then the Sanskrit for the opposing side.
\end{quote}

\subsection*{You'll want to know: \sk{विरोधः}}
\textbf{\sk{विरोधः}} — \emph{virodhaḥ} — \textquotedblleft{}opposition\textquotedblright{}.

\begin{quote}
\sk{नूतनस्य} \sk{नियमस्य} \sk{विरोधः} \sk{भवति।}
\end{quote}

\begin{quote}
\emph{nūtanasya niyamasya virodhaḥ bhavati.}
\end{quote}

\begin{quote}
There is opposition to the new rule.
\end{quote}

A masculine noun.

\begin{cousinweb}
\textbf{\sk{विरोधः}} (\emph{virodhaḥ}) is built on \emph{vi}, \textquotedblleft{}apart\textquotedblright{}, and the root √rudh, \textquotedblleft{}to obstruct\textquotedblright{}; MW gives \emph{vi-rodha} \textquotedblleft{}opposition, hostility, quarrel\textquotedblright{} (MW).
\end{cousinweb}

\begin{grammarlens}[title={What you've built}]
One more word for giving reasons and opinions.
\end{grammarlens}

\subsection*{Your turn}
\begin{itemize}
\raggedright
  \item \emph{Say it:} \emph{virodhaḥ}
  \item \emph{Say it:} \emph{virodhaḥ}, once more
  \item \emph{Say it:} \emph{nūtanasya niyamasya virodhaḥ bhavati.} — There is opposition to the new rule.
  \item \emph{Recall:} say the Sanskrit for sets out (a case), then the Sanskrit for the opposing side, then say \emph{virodhaḥ} again
\end{itemize}

\begin{tcolorbox}[breakable,colback=teal!4,colframe=teal!35!black,title={Before you move on}]
What does \sk{विरोधः} mean? (\textquotedblleft{}Opposition\textquotedblright{}.) Say it once more.
\end{tcolorbox}
\section[\texorpdfstring{\sk{समर्थयति} (samarthayati)}{समर्थयति (samarthayati)}]{\sk{समर्थयति} (samarthayati) — supports (a view)}

\label{lesson:SA-C344-samarthayati}



\begin{quote}
Before the new one: say the Sanskrit for the opposing side, then the Sanskrit for opposition.
\end{quote}

\subsection*{You'll want to know: \sk{समर्थयति}}
\textbf{\sk{समर्थयति}} — \emph{samarthayati} — \textquotedblleft{}supports (a view)\textquotedblright{}.

\begin{quote}
\sk{अहं} \sk{तव} \sk{मतं} \sk{समर्थयामि।}
\end{quote}

\begin{quote}
\emph{ahaṃ tava mataṃ samarthayāmi.}
\end{quote}

\begin{quote}
I support your view.
\end{quote}

Past (he …): \sk{सः} \sk{मम} \sk{मतं} \sk{समर्थितवान्}, \emph{saḥ mama mataṃ samarthitavān}.

\begin{cousinweb}
\textbf{\sk{समर्थयति}} (\emph{samarthayati}) is built on \emph{sam}, \textquotedblleft{}together\textquotedblright{}, and \emph{artha}, \textquotedblleft{}aim\textquotedblright{}; MW gives \emph{sam-√arth} the senses \textquotedblleft{}to make fit\textquotedblright{} and \textquotedblleft{}to judge, to consider\textquotedblright{}, and \textquotedblleft{}to support (a view)\textquotedblright{} is the sense modern Sanskrit gives it (MW).
\end{cousinweb}

\begin{grammarlens}[title={What you've built}]
One more word for giving reasons and opinions.
\end{grammarlens}

\subsection*{Your turn}
\begin{itemize}
\raggedright
  \item \emph{Say it:} \emph{samarthayati}
  \item \emph{Say it:} \emph{samarthayati}, once more
  \item \emph{Say it:} \emph{ahaṃ tava mataṃ samarthayāmi.} — I support your view.
  \item \emph{Recall:} say the Sanskrit for the opposing side, then the Sanskrit for opposition, then say \emph{samarthayati} again
\end{itemize}

\begin{tcolorbox}[breakable,colback=teal!4,colframe=teal!35!black,title={Before you move on}]
What does \sk{समर्थयति} mean? (\textquotedblleft{}Supports (a view)\textquotedblright{}.) Say it once more.
\end{tcolorbox}
\section[\texorpdfstring{\sk{अनुमन्यते} (anumanyate)}{अनुमन्यते (anumanyate)}]{\sk{अनुमन्यते} (anumanyate) — consents; approves}

\label{lesson:SA-C344-anumanyate}



\begin{quote}
Before the new one: say the Sanskrit for opposition, then the Sanskrit for supports (a view).
\end{quote}

\subsection*{You'll want to know: \sk{अनुमन्यते}}
\textbf{\sk{अनुमन्यते}} — \emph{anumanyate} — \textquotedblleft{}consents; approves\textquotedblright{}.

\begin{quote}
\sk{पिता} \sk{अनुमन्यते}, \sk{अतः} \sk{अहं} \sk{गमिष्यामि।}
\end{quote}

\begin{quote}
\emph{pitā anumanyate, ataḥ ahaṃ gamiṣyāmi.}
\end{quote}

\begin{quote}
Father consents, so I will go.
\end{quote}

Past (he …): \sk{सः} \sk{अनुमतवान्}, \emph{saḥ anumatavān}.

\begin{cousinweb}
\textbf{\sk{अनुमन्यते}} (\emph{anumanyate}) is built on \emph{anu}, \textquotedblleft{}after, along\textquotedblright{}, and the root √man, \textquotedblleft{}to think\textquotedblright{}; MW gives \emph{anu-mata} \textquotedblleft{}approved, assented to, permitted\textquotedblright{} (MW).
\end{cousinweb}

\begin{grammarlens}[title={What you've built}]
One more word for giving reasons and opinions.
\end{grammarlens}

\subsection*{Your turn}
\begin{itemize}
\raggedright
  \item \emph{Say it:} \emph{anumanyate}
  \item \emph{Say it:} \emph{anumanyate}, once more
  \item \emph{Say it:} \emph{pitā anumanyate, ataḥ ahaṃ gamiṣyāmi.} — Father consents, so I will go.
  \item \emph{Recall:} say the Sanskrit for opposition, then the Sanskrit for supports (a view), then say \emph{anumanyate} again
\end{itemize}

\begin{tcolorbox}[breakable,colback=teal!4,colframe=teal!35!black,title={Before you move on}]
What does \sk{अनुमन्यते} mean? (\textquotedblleft{}Consents; approves\textquotedblright{}.) Say it once more.
\end{tcolorbox}
\section[\texorpdfstring{\sk{ननु} (nanu)}{ननु (nanu)}]{\sk{ननु} (nanu) — is it not so; surely}

\label{lesson:SA-C344-nanu}



\begin{quote}
Before the new one: say the Sanskrit for supports (a view), then the Sanskrit for consents.
\end{quote}

\subsection*{You'll want to know: \sk{ननु}}
\textbf{\sk{ननु}} — \emph{nanu} — \textquotedblleft{}is it not so; surely\textquotedblright{}.

\begin{quote}
\sk{ननु} \sk{एतत्} \sk{उचितम्} \sk{अस्ति}?
\end{quote}

\begin{quote}
\emph{nanu etat ucitam asti?}
\end{quote}

\begin{quote}
Surely this is right, is it not?
\end{quote}

\begin{cousinweb}
\textbf{\sk{ननु}} (\emph{nanu}) is built from \emph{na}, \textquotedblleft{}not\textquotedblright{}, and \emph{nu}, \textquotedblleft{}now\textquotedblright{}; MW gives \emph{nanu} \textquotedblleft{}not?, is it not?\textquotedblright{} and so \textquotedblleft{}certainly\textquotedblright{} (MW).
\end{cousinweb}

\begin{grammarlens}[title={What you've built}]
One more word for giving reasons and opinions.
\end{grammarlens}

\subsection*{Your turn}
\begin{itemize}
\raggedright
  \item \emph{Say it:} \emph{nanu}
  \item \emph{Say it:} \emph{nanu}, once more
  \item \emph{Say it:} \emph{nanu etat ucitam asti?} — Surely this is right, is it not?
  \item \emph{Recall:} say the Sanskrit for supports (a view), then the Sanskrit for consents, then say \emph{nanu} again
\end{itemize}

\begin{tcolorbox}[breakable,colback=teal!4,colframe=teal!35!black,title={Before you move on}]
What does \sk{ननु} mean? (\textquotedblleft{}Is it not so; surely\textquotedblright{}.) Say it once more.
\end{tcolorbox}
